\section{Plan Stack}
\label{sec:plan-stack}

The Plan Stack is the persistent, ordered record of planned agent
executions for one session. It serialises the director's planning output
into a layered append-only structure that the assistant consumes in
pointer order, while permitting the director to rewrite the unexecuted
tail on failure.

\paragraph{Data model.}
The atomic unit is a \texttt{PlanStep} — one planned agent execution,
in 1:1 correspondence with one \texttt{AgentExecution} at runtime:
\begin{equation*}
\texttt{PlanStep}=(\texttt{id},\,\texttt{description},\,\texttt{status},\,\texttt{progress},\,\texttt{results}),
\end{equation*}
where $\texttt{description}\!=\!\{\texttt{agent\_id},\texttt{intent},\texttt{plan\_rationale}\}$
and $\texttt{status}\!\in\!\{\texttt{PENDING},\texttt{IN\_PROGRESS},\texttt{COMPLETED},\texttt{FAILED},\texttt{CANCELLED}\}$.
Steps are organised into ordered layers
$\mathcal{L}_0,\dots,\mathcal{L}_m$, each layer a list of
\texttt{step\_id} pointers; an \texttt{ExecutionPointer}$=(k,j)$
identifies the next step to run. The upfront plan from the director
occupies a single layer; each subsequent replan appends a new layer
that carries the rewritten tail (\S\ref{sec:director}).

\paragraph{Mutation surface.}
All structural changes go through one endpoint,
\texttt{POST /api/plan-stack/modify}, accepting an ordered list of
typed operations:
\texttt{CREATE\_STEPS}, \texttt{CREATE\_LAYERS},
\texttt{ADD\_STEPS\_TO\_LAYERS}, \texttt{REMOVE\_STEPS\_FROM\_LAYERS}.
The batch executes under a single state lock; the first failed
precondition aborts the remainder, and the response carries a partial
results list plus an error list. Two preconditions are load-bearing:
(i) steps cannot be added to a layer whose execution has already begun,
and (ii) within the active layer, insertions before the pointer are
rejected. Together they guarantee that executed history is immutable.

\paragraph{Pointer advancement.}
Given pointer $(k,j)$ and layers $\mathcal{L}_{0..m}$,
\begin{equation*}
\texttt{advance}(k,j) =
\begin{cases}
(k,\,j{+}1) & j{+}1 < |\mathcal{L}_k|, \\
(k{+}1,\,0) & j{+}1 = |\mathcal{L}_k|\;\wedge\;k{+}1\!\le\!m\;\wedge\;|\mathcal{L}_{k{+}1}|>0, \\
\bot & \text{otherwise.}
\end{cases}
\end{equation*}
Advancement is idempotent: \texttt{advance} on an exhausted pointer is
a no-op returning $\bot$, and the assistant reads the next step via the
read-only \texttt{GET /api/plan-stack/next}, so concurrent
director/executor interleaving is tolerated.


\section{Workspace}
\label{sec:workspace}

The Workspace is the per-session filesystem and in-memory state shared
by every sub-agent execution within a session. It is the only writable
surface below the assistant: agents do not perform direct I/O, and
artefacts are addressable purely through the workspace caption index.

\paragraph{Layout.}
A workspace rooted at \texttt{Runtime/\{workspace\_id\}/} contains
\texttt{global\_memory.md} (a Markdown-wrapped JSON array of
\texttt{GlobalMemoryEntry} records), \texttt{logs.jsonl}, an
\texttt{artifacts/\{agent\_id\}/} tree of binaries and JSON snapshots,
and an \texttt{inputs/} tree of raw user uploads. The
\texttt{Workspace} object composes four collaborators —
\texttt{FileManager}, \texttt{LogManager}, \texttt{GlobalMemory},
\texttt{ArtifactWriter} — each owning one slice of this state.

\paragraph{Cross-agent data contract.}
Producers register an \texttt{ArtifactRef}$=(\texttt{caption},\texttt{scope},\texttt{path},\texttt{mime})$;
consumers receive a \texttt{ResolvedArtifactEntry} with the same four
fields plus an optional \texttt{payload} dict (loaded only for
\texttt{application/json}). Every cross-agent data channel passes
through this single shape — \texttt{(single)} labels resolve to one
entry, \texttt{(collection)} labels resolve to a list. Consumers never
parse producer-internal field names; the canonical idiom is to JSON-dump
\texttt{entry.payload} into the consumer's prompt and let its LLM read
the upstream structure as text.\footnote{This is the project-wide
``strict-out / loose-in'' convention: producers maintain their own
output schema rigorously (Pydantic + evaluator), consumers make zero
structural assumptions about it.}

\paragraph{Caption index and InputResolver.}
\texttt{GlobalMemory} indexes every workspace artefact by
\emph{caption} — a short natural-language sentence describing the
artefact's role and producer. Given a consumer agent and its
\texttt{input\_needs\_description} (a list of
\texttt{[label] (single|collection): nl-description} blocks),
\texttt{InputResolver} (i) renders the caption index and label specs
into a single LLM prompt, (ii) parses the LLM's
$\{\texttt{label}\!\to\![\texttt{ids}]\}$ selection, (iii) maps each
chosen id back to a path, and (iv) materialises a
\texttt{ResolvedArtifactEntry} per path (loading JSON payload on
demand). The output is the
$\texttt{resolved\_artifacts}\!:\!\{\texttt{label}\!\to\!\texttt{Entry}\,|\,[\texttt{Entry}]\}$
dict that the assistant hands to \texttt{descriptor.build\_input}
(\S\ref{sec:assistant}).

\paragraph{Persistence.}
After a successful execution, the assistant calls
\texttt{workspace.persist\_execution\_from\_plan(execution, base\_plan,
overwrite\_existing, captions)}. \texttt{ArtifactWriter} writes the JSON
snapshot and binary media into the workspace tree, and registers a new
\texttt{GlobalMemoryEntry} with one \texttt{ArtifactRef} per artefact.
With \texttt{overwrite\_existing=True} the writer first prunes prior
artefacts produced by the same \texttt{agent\_id}: physical files are
unlinked and registry entries removed, so re-execution produces a clean
replacement rather than an accumulating duplicate set. Captions are
authored by the descriptor (one per artefact); they describe the
artefact's role, never its content.\footnote{Inlining content into a
caption gives no recall lift but creates near-duplicate captions that
collide in the InputResolver LLM's selection.}


\section{Assistant}
\label{sec:assistant}

The Assistant is the stateless executor of one (\texttt{agent\_id},
\texttt{step\_id}) pair against one \texttt{Workspace}. It owns the
descriptor-driven dispatch pipeline; it does not own planning,
sequencing, or failure recovery (those belong to the director —
\S\ref{sec:director}).

\paragraph{Entry point.}
A single HTTP endpoint
\texttt{POST /api/assistant/execute} with body
$\{\texttt{agent\_id},\texttt{step\_id}\}$ invokes
\begin{equation*}
\texttt{execute\_agent\_for\_step}(\texttt{agent\_id},\texttt{step\_id})\!\to\!\texttt{AgentExecution},
\end{equation*}
where the returned
$\texttt{AgentExecution}=(\texttt{id},\texttt{agent\_id},\texttt{step\_id},\texttt{status},\texttt{inputs},\texttt{results},\texttt{error},\texttt{resolved\_input\_paths})$
records both the resolved input set and the post-execution outputs.

\paragraph{Three-boundary pipeline.}
\begin{enumerate}
\item \texttt{build\_execution\_inputs}: validate \texttt{agent\_id} in
the registry, run \texttt{InputResolver} against the workspace caption
index using the descriptor's \texttt{input\_needs\_description}, return
$\{\texttt{step\_id},\texttt{resolved\_artifacts},\texttt{resolved\_input\_paths}\}$.
\item \texttt{execute\_agent}: call
\texttt{descriptor.build\_equipped\_agent(llm)} (which wires the
agent's evaluator and optional materializer), then
\texttt{descriptor.build\_input(step\_id, resolved\_artifacts)} to
produce the typed Pydantic input, then \texttt{agent.run(typed\_input,
materialize\_ctx)} which performs LLM generation, three-tier evaluation
(structural / creative / asset), and media materialisation.
\item \texttt{process\_results}: derive a deterministic persist plan
from \texttt{execution.results}, call
\texttt{workspace.persist\_execution\_from\_plan} with descriptor-authored
captions, optionally promote consumed raw inputs to global scope.
\end{enumerate}

\paragraph{Quality gate.}
The agent's three-tier evaluator returns a Boolean
$\texttt{passed}\!\in\!\{0,1\}$. The assistant maps
\begin{equation*}
\texttt{ExecutionStatus}=
\begin{cases}
\texttt{COMPLETED} & \texttt{passed}=1, \\
\texttt{FAILED} & \texttt{passed}=0,
\end{cases}
\end{equation*}
and stamps \texttt{error} from a structured debug record. Two failure
classes are distinguished by prefix in the error string:
$[\texttt{upstream\_input\_rejected}]$ means the agent rejected its
resolved inputs (a producer-side problem), while plain
\texttt{quality gate failed: \dots} means the agent's own output failed
its evaluator. The director's replanner reads the prefix to decide
whether to retry the same step or rewrite an upstream tail
(\S\ref{sec:director}).

\paragraph{Descriptor as the only extension point.}
A new agent is added by registering a
\texttt{SubAgentDescriptor}$=$(\texttt{agent\_id},
\texttt{catalog\_entry}, \texttt{agent\_factory},
\texttt{evaluator\_factory}, \texttt{build\_input},
\texttt{build\_captions}, \texttt{materializer\_factory},
\texttt{input\_needs\_description}, \texttt{promotes\_consumed\_inputs\_to\_global})
in \texttt{AGENT\_REGISTRY}. The assistant contains zero
agent-specific code paths: every divergence between agents lives in the
descriptor.


\section{Director}
\label{sec:director}

The Director is the stateless agent that converts a user goal into a
chain of plan steps and reacts to runtime failures. It is the only
producer of \texttt{PlanStep.description} content; it never executes
agents itself.

\paragraph{Upfront planning.}
Given a merged session goal $g$ and the registered agent set
$\mathcal{A}$, the planner $\pi_\theta$ samples
\begin{equation*}
\pi_\theta(g,\mathcal{A})\!\to\!\{\texttt{rationale},\,\hat{\mathbf{c}}=(\hat{c}_1,\dots,\hat{c}_n)\},\quad \hat{c}_i=(\texttt{agent\_id}_i,\texttt{intent}_i),
\end{equation*}
with $\hat{c}_i.\texttt{agent\_id}\!\in\!\mathcal{A}$, all
\texttt{agent\_id}s pairwise distinct, and
$n\!\le\!\texttt{MAX\_PIPELINE\_STEPS}\!=\!20$. The system prompt is the
canonical \texttt{PLAN\_UPFRONT\_CORE} (\S\ref{sec:sft}) — task
definition, output schema, structural rules, the allowed-id list, and
a per-agent descriptor block for every $a\!\in\!\mathcal{A}$ — with
optional scaffold flags (\texttt{fewshots}, \texttt{policies},
\texttt{topology}) toggling additional content. Plan validation
discards entries with unregistered \texttt{agent\_id}, enforces the
length cap, and rejects empty plans; on validation failure the planner
is re-sampled up to $3$ times.

\paragraph{Persistence.}
The accepted chain is committed in two atomic batches against
\texttt{POST /api/plan-stack/modify}: (i) \texttt{CREATE\_STEPS} with
each step's $\texttt{description}=\{\texttt{agent\_id},\texttt{intent},\texttt{plan\_rationale}\}$,
plus \texttt{CREATE\_LAYERS} for one new layer at index $m{+}1$;
(ii) \texttt{ADD\_STEPS\_TO\_LAYERS} mapping the freshly created
\texttt{step\_id}s into that layer in chain order. The execution
pointer is initialised to the new layer's first slot if previously
null.

\paragraph{Execution loop and failure handling.}
The director then runs an executor loop bounded by
\texttt{MAX\_EXECUTIONS\_PER\_CYCLE}$=30$:
\begin{equation*}
\begin{array}{l}
\texttt{while budget>0:}\\
\quad s\!=\!\texttt{get\_next\_step()};\;\;\texttt{if not}\;s\colon\;\texttt{return}\\
\quad \texttt{set\_status}(s,\texttt{IN\_PROGRESS})\\
\quad e\!=\!\texttt{execute\_agent}(s.\texttt{agent\_id},s.\texttt{id})\\
\quad \texttt{set\_status}(s,e.\texttt{status})\\
\quad \texttt{if}\;e.\texttt{status}=\texttt{FAILED}\colon\;\texttt{handle\_failure}(s,e)\\
\quad \texttt{advance\_execution\_pointer()}.
\end{array}
\end{equation*}
On \texttt{FAILED}, \texttt{handle\_failure} consults a
\texttt{ReplanDecision} from
\texttt{planner.replan\_on\_failure(failed\_step, pending\_tail,
completed\_tail, user\_goal)}. Three branches follow:
\begin{itemize}
\item \emph{retry} — same \texttt{agent\_id} re-emitted as the head of
\texttt{new\_tail}; director resets the failed step to \texttt{PENDING}
and rolls the pointer back so it re-runs.
\item \emph{replan} — \texttt{new\_tail} is non-empty and differs from
the pending tail; director removes all \texttt{PENDING} steps from the
stack via \texttt{REMOVE\_STEPS\_FROM\_LAYERS} and appends a fresh
layer carrying \texttt{new\_tail}; the pointer naturally advances into
that layer.
\item \emph{skip} — \texttt{new\_tail} is empty; the stack is
unchanged, the pointer simply moves past the failed step on the next
\texttt{advance}.
\end{itemize}
The total number of \emph{replan} or \emph{retry} actions per user turn
is capped by $\texttt{MAX\_REPLAN\_ROUNDS}\!=\!2$; once exhausted, the
director runs the remaining stack without further rewriting.

\paragraph{Memory projection.}
Both the upfront and replan prompts receive a \emph{slim} projection
of the current stack (last $\texttt{DIRECTOR\_MEMORY\_WINDOW}\!=\!20$
steps): per step the planner sees \texttt{step\_id},
\texttt{agent\_id}, \texttt{status}, \texttt{intent}, and a
results-summary truncated to $400$ characters. Failure messages in the
replan prompt are truncated to $2000$ characters. The full workspace
state — captions, payloads, binaries — is never surfaced to the
director; routing decisions are taken on metadata alone.\footnote{This
is what makes director routing trainable as a standalone task
(\S\ref{sec:bench}): the inputs to $\pi_\theta$ are textual and
bounded, and the output is a closed-vocabulary chain.}
